\documentclass[11pt,a4paper]{article}

% --- {{fr:Encodage et langue|en:Encoding and language}} ---
\usepackage[T1]{fontenc}
\usepackage[{{babel}}]{babel}
\usepackage{lmodern}
\usepackage{microtype}

% --- {{fr:Mise en page|en:Page layout}} ---
\usepackage[margin=2.5cm]{geometry}

% --- {{fr:Mathématiques|en:Mathematics}} ---
\usepackage{amsmath, amssymb, amsthm}

% --- {{fr:Figures et tableaux|en:Figures and tables}} ---
\usepackage{graphicx}
\usepackage{booktabs}
\usepackage{xcolor}

% --- {{fr:Liens (à charger en dernier)|en:Links (load last)}} ---
\usepackage[hidelinks]{hyperref}

\title{{{title}}}
\author{{{author}}}
\date{{{date}}}

\begin{document}

\maketitle

\begin{abstract}
  {{fr:Résumez ici votre article en quelques lignes.|en:Summarise your article here in a few lines.}}
\end{abstract}

\section{Introduction}
\label{sec:introduction}

{{fr:Voici un paragraphe. Laissez une ligne vide pour commencer un nouveau paragraphe.|en:Here is a paragraph. Leave a blank line to start a new paragraph.}}
{{fr:Les formules en ligne s'écrivent entre dollars : $e^{i\pi} + 1 = 0$.|en:Inline formulas go between dollars: $e^{i\pi} + 1 = 0$.}}

\section{{{fr:Résultats|en:Results}}}
\label{sec:resultats}

{{fr:L'équation~\eqref{eq:gauss} est numérotée et peut être citée.|en:Equation~\eqref{eq:gauss} is numbered and can be referenced.}}
\begin{equation}
  \int_{-\infty}^{+\infty} e^{-x^2} \, \mathrm{d}x = \sqrt{\pi}
  \label{eq:gauss}
\end{equation}

{{fr:La figure~\ref{fig:exemple} et le tableau~\ref{tab:exemple} illustrent les résultats.|en:Figure~\ref{fig:exemple} and Table~\ref{tab:exemple} illustrate the results.}}

\begin{figure}[htbp]
  \centering
  % \includegraphics[width=0.7\linewidth]{figure}
  \fbox{\parbox{0.6\linewidth}{\centering\vspace{2cm}{{fr:Placez votre image ici|en:Put your image here}}\vspace{2cm}}}
  \caption{{{fr:Une figure d'exemple.|en:An example figure.}}}
  \label{fig:exemple}
\end{figure}

\begin{table}[htbp]
  \centering
  \caption{{{fr:Un tableau d'exemple.|en:An example table.}}}
  \label{tab:exemple}
  \begin{tabular}{lcc}
    \toprule
    {{fr:Méthode|en:Method}} & {{fr:Temps (s)|en:Time (s)}} & {{fr:Précision|en:Accuracy}} \\
    \midrule
    A & 1{,}2 & 0{,}91 \\
    B & 0{,}8 & 0{,}87 \\
    \bottomrule
  \end{tabular}
\end{table}

\section{Conclusion}

{{fr:Voir la section~\ref{sec:introduction}.|en:See Section~\ref{sec:introduction}.}}

\end{document}
